\documentclass[12pt]{article}
\usepackage{ctex}
\usepackage{amsmath, amssymb}
\usepackage{geometry}
\geometry{a4paper, margin=2.5cm}
\usepackage{enumitem}
\title{第 3 章\ 功和能\ 例题}
\author{}
\date{}

\begin{document}
\maketitle

\begin{enumerate}[label=\textbf{例\arabic*.}, leftmargin=*]

% 例 1
\item 已知 $m = 2\,\mathrm{kg}$，在 $F = 12t$（N）作用下由静止做直线运动。求 $t = 0 \to 2\,\mathrm{s}$ 内 $F$ 做的功及 $t = 2\,\mathrm{s}$ 时的瞬时功率。

% 例 2
\item 质量 $m = 10\,\mathrm{kg}$ 的质点在外力作用下做平面曲线运动，速度 $\vec{v} = 4t^2\vec{i} + 16\vec{j}$，开始时质点位于坐标原点。求质点从 $y = 16\,\mathrm{m}$ 到 $y = 32\,\mathrm{m}$ 过程中外力做的功。

% 例 3
\item 用力 $\vec{F}$ 缓慢拉质量为 $m$ 的小球，$\vec{F}$ 保持方向不变。求 $\theta = \theta_0$ 时 $\vec{F}$ 所做的功。

% 例 4
\item 长 $L$、质量 $M$ 的均匀柔绳，$A$ 端挂在天花板上自然下垂，将 $B$ 端沿铅直方向提高到与 $A$ 端同高处。求该过程中重力所做的功。

% 例 5
\item 蓄水池面积 $S$、水深 $h$、水面距地面 $H$、水的密度为 $\rho$。求将水全部抽出需要做的功。

% 例 6
\item 风力 $F$ 作用于向北运动的船，风力方向变化规律 $\theta = Bs$，其中 $s$ 为位移、$B$ 为常数、$\theta$ 为 $F$ 与 $S$ 间的夹角。若风的方向自南变到东，求风力做的功。

% 例 7
\item 思考题：作用于质点的力 $\vec{F} = y\vec{i} + 3x^2\vec{j}$（N），求该力沿 $oab$ 与 $ob$ 路径所做的功（$a(2,0), b(2,1)$）。

% 例 8
\item 物体质量 $m$，弹簧劲度系数 $k$，$A$ 板及弹簧质量不计。求自弹簧原长 $O$ 处突然无初速度地加上物体时，弹簧的最大压缩量 $\lambda_{\max}$。

% 例 9
\item 长为 $l$ 的均质链条，部分置于水平面上，另一部分自然下垂，链条与水平面间静摩擦系数 $\mu_0$、滑动摩擦系数 $\mu$。\\
（1）满足什么条件链条将开始滑动？\\
（2）若下垂部分长度为 $b$ 时链条自静止开始滑动，当链条末端刚刚滑离桌面时其速度为多少？

% 例 10
\item 水平面内有半径为 $R$ 的圆，离圆心 $O$ 距离 $S$ 处有质量很大、视为固定的力心 $O'$，力心对单位质量的有心引力为 $\mu r$（$r$ 为力心到质点 $Q$ 的位矢大小），质点 $Q$ 被限制在圆周上运动。求：\\
（1）质点 $Q$ 从 $B$ 点由静止出发转过 $\varphi$ 角有心力做的功；\\
（2）质点通过第二象限所经历的时间。

% 例 11
\item 物体质量 $m$、弹簧劲度系数 $k$，自弹簧原长无初速度地放上物体。求弹簧的最大压缩量 $y_{\max}$。

% 例 12
\item 用弹簧连接两木板 $m_1$、$m_2$，弹簧压缩 $x_0$。求给 $m_2$ 加多大的压力能使 $m_1$ 离开桌面。

% 例 13
\item 物体 $M$ 悬于弹簧上，弹簧弹性系数 $k$，弹簧原长等于圆环半径，不计摩擦。求物体自弹簧原长无初速度地沿圆环滑至最低点 $B$ 时所获得的动能。

% 例 14
\item 物体以第二宇宙速度 $v_0 = \sqrt{2GM_e/R_e}$ 从地球表面沿铅直方向发射，阻力忽略。求物体从地面飞行到与地心相距 $nR_e$ 处所经历的时间。

\end{enumerate}

\end{document}
